\begin{abstract}
Novel view synthesis has evolved rapidly, advancing from Neural Radiance Fields to 3D Gaussian Splatting (3DGS), which offers real-time rendering and rapid training without compromising visual fidelity. However, 3DGS relies heavily on accurate camera poses and high-quality point cloud initialization, which are difficult to obtain in sparse-view scenarios. While traditional Structure from Motion (SfM) pipelines often fail in these settings, existing learning-based point estimation alternatives typically require reliable reference views and remain sensitive to pose or depth errors. In this work, we propose a robust method utilizing $\pi^3$, a reference-free point cloud estimation network. We integrate dense initialization from $\pi^3$ with a regularization scheme designed to mitigate geometric inaccuracies. Specifically, we employ uncertainty-guided depth supervision, normal consistency loss, and depth warping. Experimental results demonstrate that our approach achieves state-of-the-art performance on the Tanks and Temples, LLFF, DTU, and MipNeRF360 datasets.

%Novel view synthesis has seen significant advancements, starting with traditional Structure from Motion (SfM) pipelines and advancing to Neural Radiance Fields and 3D Gaussian Splatting (3DGS), which combines minimal quality trade-offs with real-time rendering and substantially faster training. 3DGS relies on accurate camera poses and an initial point cloud, which is challenging in sparse-view settings. Traditional SfM often struggles in such sparse-view settings, which gave rise to point cloud estimation networks. These estimation networks require a reliable reference view and struggle with inaccurate depth estimations and camera poses. In this thesis, we propose a new method that leverages $\pi^3$ as a point cloud estimation network, which does not require reference views. Our method combines the dense initialization from the inferred point cloud with optimized regularization to handle inaccuracies in $\pi^3$ depth maps. These regularizations include uncertainty-guided depth loss, normal loss and depth warping. Our model achieves state-of-the-art results across different datasets, such as Tanks and Temples, LLFF, DTU and MipNeRF360.
\end{abstract}